\subsection{Linear Systems and Stability}
% Matrix exponentials, stable/unstable subspaces, Lyapunov stability, hyperbolicity.

Scalar first-order linear equations admit an explicit evolution factor and a variation-of-constants representation.

\begin{proposition} \label{prop:scalar-linear-variation-of-constants} \hlt{Scalar Variation of Constants}\\
Let $I\subseteq\R$ be open, $a,g\in C(I,\R)$, and $(t_0,x_0)\in I\times\R$.\\
Put $E(t,s)\defeq\exp\left(\int_s^t a(r)\,dr\right)$ and $\phi(t)\defeq E(t,t_0)x_0+\int_{t_0}^tE(t,s)g(s)\,ds$. Then:
\begin{enumerate}[label=(\roman*)]
\setlength{\itemsep}{0pt}
\item $\phi\in C^1(I,\R)$ and $(\phi'-a\phi,\phi(t_0))=(g,x_0)$.
\item If $x\in C^1(I,\R)$ and $(x'-ax,x(t_0))=(g,x_0)$, then $x=\phi$.
\end{enumerate}
\begin{proof}
\begin{enumerate}[label=(\roman*)]
\setlength{\itemsep}{0pt}
\item Let $B(t)\defeq x_0+\int_{t_0}^tE(t_0,s)g(s)\,ds$. By additivity, $E(t,s)=E(t,t_0)E(t_0,s)$, so $\phi(t)=E(t,t_0)B(t)$.\\
By Theorem \ref{thm:integral-function-continuity-and-derivative}, $B\in C^1$, $\partial_tE(t,t_0)=a(t)E(t,t_0)$, and $B'(t)=E(t_0,t)g(t)$.\\
Hence $\phi\in C^1$, $\phi'(t)=a(t)\phi(t)+E(t,t_0)E(t_0,t)g(t)=a(t)\phi(t)+g(t)$, and $\phi(t_0)=x_0$.
\item If $(x'-ax,x(t_0))=(g,x_0)$, then $\frac{d}{dt}(E(t_0,t)x(t))=E(t_0,t)(x'(t)-a(t)x(t))=E(t_0,t)g(t)$.\\
By Theorem \ref{thm:fundamental-theorem-calculus-riemann}, $E(t_0,t)x(t)=x_0+\int_{t_0}^tE(t_0,s)g(s)\,ds$.\\
Multiplying by $E(t,t_0)$ and using $E(t,t_0)E(t_0,s)=E(t,s)$ gives $x(t)=\phi(t)$.
\end{enumerate}
\hlt{Structural Commentary}: $E(t_0,t)$ removes $a(t)x$; integration accumulates $g$.\\
\hlt{Key Insight}: $(E(t_0,t)x)'=E(t_0,t)g$.
\end{proof}
\end{proposition}

\begin{proposition} \label{prop:fiber-preserving-ode-transformation} \hlt{Fiber-Preserving Change of Variables}\\
Let $\Omega,\widetilde{\Omega}\subseteq\R^2$ be open, $f:\Omega\to\R$, and $T:\Omega\to\widetilde{\Omega}$ be a $C^1$ diffeomorphism.\\
Assume $T(t,x)=(\sigma(t),\eta(t,x))$ and $T^{-1}(s,y)=(\tau(s),\xi(s,y))$. Put $G\defeq\partial_t\eta+\partial_x\eta\,f$.\\
Define $\widetilde f:\widetilde\Omega\to\R$ by $\widetilde f(s,y)\defeq\tau'(s)G(\tau(s),\xi(s,y))$.\\
Let $I\subseteq\R$ be an interval, $x\in C^1(I,\R)$, $\Gamma_x\defeq\{(t,x(t)):t\in I\}\subseteq\Omega$, and $J\defeq\sigma(I)$.\\
Set $y:J\to\R$ by $y(s)\defeq\eta(\tau(s),x(\tau(s)))$. Then $y\in C^1(J,\R)$, and $x'=f(\,\cdot\,,x)$ on $I$ iff $y'=\widetilde f(\,\cdot\,,y)$ on $J$.
\begin{proof}
Since $\det DT(t,x)=\sigma'(t)\partial_x\eta(t,x)\neq0$, $\sigma'\neq0$, $\partial_x\eta\neq0$, and $\tau'(\sigma(t))\sigma'(t)=1$.
\begin{enumerate}[label=\roman*.]
\setlength{\itemsep}{0pt}
\item[$(\Rightarrow)$] If $x'=f(\,\cdot\,,x)$ and $s\in J$, then $y'(s)=\tau'(s)G(\tau(s),x(\tau(s)))$. Since $\xi(s,y(s))=x(\tau(s))$, $y'(s)=\widetilde f(s,y(s))$.
\item[$(\Leftarrow)$] Let $t\in I$. Differentiating $y(\sigma(t))=\eta(t,x(t))$ gives $\sigma'(t)y'(\sigma(t))=\partial_t\eta(t,x(t))+\partial_x\eta(t,x(t))x'(t)$.\\
Since $y'=\widetilde f(\,\cdot\,,y)$, $\sigma'(t)y'(\sigma(t))=G(t,x(t))$. Hence $\partial_x\eta(t,x(t))[x'(t)-f(t,x(t))]=0$, so $x'=f(\,\cdot\,,x)$.
\end{enumerate}
\hlt{Structural Commentary}: Transform graphs; apply the chain rule.\\
\hlt{Key Insight}: $\widetilde f$ is the pushed-forward vector field.
\end{proof}
\end{proposition}

\begin{method} \label{meth:homogeneous-first-order-reduction} \hlt{Homogeneous First-Order Reduction}\\
For $x'=f(x/t)$ on an interval with $t\neq0$:
\begin{enumerate}[label=(\roman*)]
\setlength{\itemsep}{0pt}
\item Set $y\defeq x/t$, so $x=ty$.
\item Differentiate to obtain $y'=(tx'-x)/t^2=[f(y)-y]/t$.
\item Each root $c$ of $f(c)=c$ gives the constant-ratio solution $x(t)=ct$.
\item On every branch where $f(y)\neq y$, separate $dy/(f(y)-y)=dt/t$, solve for $y$, and recover $x=ty$.
\end{enumerate}
\end{method}

\begin{method} \label{meth:affine-fractional-first-order-reduction} \hlt{Affine-Fractional First-Order Reduction}\\
Let $A(t,x)\defeq ax+bt+c$, $B(t,x)\defeq\alpha x+\beta t+\gamma$, $B\neq0$, and $\Delta\defeq a\beta-\alpha b$. For $x'=f(A/B)$:
\begin{enumerate}[label=(\roman*)]
\setlength{\itemsep}{0pt}
\item If $\Delta\neq0$, choose $(t_*,x_*)$ with $A(t_*,x_*)=B(t_*,x_*)=0$. Set $s\defeq t-t_*$ and $y\defeq x-x_*$.\\
Then $\frac{dy}{ds}=f(\frac{ay+bs}{\alpha y+\beta s})$; for $s\neq0$, $z\defeq y/s$ satisfies $s\frac{dz}{ds}=f(\frac{az+b}{\alpha z+\beta})-z$.
\item If $\Delta=0$, $(a,b)\neq(0,0)$, and $(\alpha,\beta)=\lambda(a,b)$, set $\widehat{f}(u)\defeq f(\frac{u+c}{\lambda u+\gamma})$.\\
If $a\neq0$, use $u\defeq ax+bt$: $u'=a\widehat{f}(u)+b$, with $x=\frac{u-bt}{a}$. If $a=0$, then $x'=\widehat{f}(bt)$.
\end{enumerate}
\end{method}

\begin{example} \label{ex:euler-equation-log-time} \hlt{Euler Equation Under Logarithmic Time}\\
For $t>0$, set $s\defeq\log t$ and $y(s)\defeq x(e^s)$. Then $x'(t)=t^{-1}y'(s)$ and $x''(t)=t^{-2}(y''(s)-y'(s))$.\\
Hence Euler equation $t^2x''+3tx'+x=2/t$ becomes constant-coefficient equation $y''+2y'+y=2e^{-s}$.
\end{example}

\begin{proposition} \label{prop:exact-first-order-equation-potential} \hlt{Exact First-Order Equations}\\
Let $R\defeq I\times J$ be an open rectangle, let $p,q\in C^1(R,\R)$, and fix $(x_0,y_0)\in R$. The following are equivalent:
\begin{enumerate}[label=(\roman*)]
\setlength{\itemsep}{0pt}
\item $\exists F\in C^2(R,\R)$ with $\partial_yF=p$ and $\partial_xF=q$.
\item $\partial_xp=\partial_yq$ on $R$.
\end{enumerate}
Assume (i)--(ii), and set $F_0(x,y)\defeq\int_{x_0}^{x}q(s,y_0)\,ds+\int_{y_0}^{y}p(x,r)\,dr$.\\
Then $F_0$ satisfies (i). If $F_1$ satisfies (i), then $\exists c\in\R$ with $F_1=F_0+c$.\\
Moreover, for every interval $L\subseteq I$ and $y\in C^1(L,J)$, $p(\,\cdot\,,y)y'+q(\,\cdot\,,y)=0$ iff $F_0(\,\cdot\,,y)$ is constant on $L$.
\begin{proof}
\begin{enumerate}[label=(\roman*)]
\setlength{\itemsep}{0pt}
\item[$(i)$]$\Rightarrow (ii)$ If $\partial_yF=p$ and $\partial_xF=q$, then by Corollary \ref{cor:mixed-partials-from-second-derivative}, $\partial_xp=\partial_x\partial_yF=\partial_y\partial_xF=\partial_yq$.
\item[$(ii)$]$\Rightarrow (i)$ For $F_0$, $\partial_yF_0=p$ and $\partial_xF_0=q(x,y_0)+\int_{y_0}^{y}\partial_xp(x,r)\,dr$.\\
By (ii), $\partial_xF_0=q(x,y_0)+\int_{y_0}^{y}\partial_yq(x,r)\,dr=q(x,y)$. Since $p,q\in C^1(R,\R)$, $F_0\in C^2(R,\R)$.
\end{enumerate}
If $F_1$ also satisfies (i), set $H\defeq F_1-F_0$. Then $\nabla H=0$.\\
By the one-variable fundamental theorem along rectangle sides, $H(x,y)=H(x_0,y_0)$; hence $F_1=F_0+c$.\\
For $y\in C^1(L,J)$, the chain rule gives $(F_0(\,\cdot\,,y))'=q(\,\cdot\,,y)+p(\,\cdot\,,y)y'$.\\
Hence $p(\,\cdot\,,y)y'+q(\,\cdot\,,y)=0$ iff $F_0(\,\cdot\,,y)$ is constant on $L$.\\
\hlt{Structural Commentary}: Build a potential; differentiate along graphs.\\
\hlt{Key Insight}: Mixed partials are the compatibility condition.
\end{proof}
\end{proposition}

\begin{method} \label{meth:integrating-factor-dependent-on-x} \hlt{Integrating Factor Depending Only on $x$}\\
Let $R\subseteq\R^2$ be a rectangle, let $p,q\in C^1(R,\R)$, and assume $p\neq0$ on $R$. For $p(x,y)y'+q(x,y)=0$:
\begin{enumerate}[label=(\roman*)]
\setlength{\itemsep}{0pt}
\item Set $r(x,y)\defeq(\partial_yq-\partial_xp)/p$.
\item If $r(x,y)=r_0(x)$ on $R$, set $\mu(x)\defeq\exp(\int_{x_0}^{x}r_0(s)\,ds)$.
\item Multiply by $\mu$: $\partial_x(\mu p)=\mu(r_0p+\partial_xp)=\mu\partial_yq=\partial_y(\mu q)$, so $\mu p\,y'+\mu q=0$ is exact by Proposition \ref{prop:exact-first-order-equation-potential}.
\end{enumerate}
For example, $xy'+3x-2y=0$ has $r_0(x)=-3/x$ on $x\neq0$, so one may take $\mu(x)=x^{-3}$.\\
A potential is $F(x,y)=y/x^2-3/x$, and the solutions are $y=3x+Cx^2$.
\end{method}

\begin{method} \label{meth:bernoulli-equation-linearization} \hlt{Bernoulli Linearization}\\
Let $I\subseteq\R$ be open, $a,b\in C(I,\R)$, $n\in\R\setminus\{0,1\}$. Restrict to subinterval $J\subseteq I$, $x>0$ for $x'=a(t)x+b(t)x^n$.
\begin{enumerate}[label=(\roman*)]
\setlength{\itemsep}{0pt}
\item Set $y\defeq x^{1-n}$.
\item Linearize: $y'=(1-n)x^{-n}x'=(1-n)a(t)y+(1-n)b(t)$ on $J$.
\item Solve $y'=\alpha y+\beta$ with $\alpha\defeq(1-n)a$, $\beta\defeq(1-n)b$ (Proposition \ref{prop:scalar-linear-variation-of-constants}); recover $x=y^{1/(1-n)}$ on $\{y>0\}$.
\end{enumerate}
\end{method}

\begin{method} \label{meth:riccati-equation-linearization} \hlt{Riccati Linearization from a Particular Solution}\\
Let $I\subseteq\R$ be open, $a,b,c\in C(I,\R)$, and let $x_p\in C^1(I,\R)$ solve $x_p'=ax_p+bx_p^2+c$.\\
Work on a subinterval $J\subseteq I$ with $x\neq x_p$ for $x'=ax+bx^2+c$.
\begin{enumerate}[label=(\roman*)]
\setlength{\itemsep}{0pt}
\item Set $u\defeq x-x_p$ and $y\defeq u^{-1}$ on $J$.
\item Subtract and invert: $u'=(a+2bx_p)u+bu^2$, so $y'=-(a+2bx_p)y-b$.
\item Solve $y'=\alpha y+\beta$ with $\alpha\defeq-(a+2bx_p)$, $\beta\defeq-b$ (Proposition \ref{prop:scalar-linear-variation-of-constants}); recover $x=x_p+y^{-1}$ on $\{y\neq0\}$.
\end{enumerate}
\end{method}

\begin{method} \label{meth:solve-linear-ode-by-jordan-form} \hlt{Solving $\mathbf{x}'(t)=A\mathbf{x}(t)$ by Jordan Form}\\
Let $A\in\C^{n\times n}$ and $\mathbf{x}_0\in\C^n$. To solve $\mathbf{x}'(t)=A\mathbf{x}(t)$ with $\mathbf{x}(0)=\mathbf{x}_0$:
\begin{enumerate}[label=(\roman*)]
\setlength{\itemsep}{0pt}
\item By Theorem \ref{thm:complex-jordan-form}, choose $S\in GL_n(\C)$ with $A=SJS^{-1}$, where $J=\bigoplus\limits_{i=1}^rJ_{k_i}(\lambda_i)$.
\item Set $\mathbf{y}(t)=S^{-1}\mathbf{x}(t)$ and $\mathbf{y}_0=S^{-1}\mathbf{x}_0$. Then $\mathbf{y}'(t)=J\mathbf{y}(t)$ and $\mathbf{x}(t)=S\mathbf{y}(t)$.
\item Split $\mathbf{y}=(\mathbf{y}_1,\ldots,\mathbf{y}_r)$ according to the Jordan blocks. Then $\mathbf{y}_i'(t)=J_{k_i}(\lambda_i)\mathbf{y}_i(t)$.
\item By Proposition \ref{prop:exponential-of-jordan-block}, $\mathbf{y}_i(t)=e^{tJ_{k_i}(\lambda_i)}\mathbf{y}_i(0)=e^{\lambda_i t}\sum\limits_{q=0}^{k_i-1}\frac{t^q}{q!}J_{k_i}(0)^q\mathbf{y}_i(0)$ for $1\leq i\leq r$.\\
Differentiating verifies the block equation.
\item Reassemble $\mathbf{y}(t)$ and transform back: $\mathbf{x}(t)=S\mathbf{y}(t)$.
\end{enumerate}
Thus a size-$k$ block with eigenvalue $\lambda$ contributes entries of the form $e^{\lambda t}p(t)$, where $\deg p<k$.
\end{method}

\begin{perspective} \label{per:jordan-functions-polynomial-factors} \hlt{Jordan Functions and Polynomial Factors}\\
For a Jordan block $J_k(\lambda)=\lambda I_k+N$ with $N^k=0$, any polynomial $f$ satisfies $f(J_k(\lambda))=\sum\limits_{r=0}^{k-1}\frac{f^{(r)}(\lambda)}{r!}N^r$.\\
The nilpotent part contributes derivative terms; for the exponential these become polynomial factors in $t$ by Proposition \ref{prop:exponential-of-jordan-block}. As $t\to+\infty$, $\operatorname{Re}\lambda<0$ gives blockwise decay and $\operatorname{Re}\lambda>0$ gives growth.\\
If $\operatorname{Re}\lambda=0$, the block evolution is bounded iff $N=0$.
\end{perspective}

\begin{example} \label{ex:nilpotent-jordan-block-linear-evolution} \hlt{Nilpotent Jordan Block Evolution}\\
Let $N\defeq J_3(0)=\left(\begin{smallmatrix}0&1&0\\0&0&1\\0&0&0\end{smallmatrix}\right)$, so $N^3=0$. By Proposition \ref{prop:exponential-of-jordan-block}, $e^{tN}=I_3+tN+\frac{t^2}{2}N^2=\left(\begin{smallmatrix}1&t&\frac{t^2}{2}\\0&1&t\\0&0&1\end{smallmatrix}\right)$. Nilpotent Jordan blocks create polynomial, rather than exponential, growth along generalized-eigenvector chains.
\end{example}
